\section{Small ideal kernels at every frequency}\label{sec:ideal}

We use the groups and probability measures of \Cref{sec:transference}.
In this section each big group $\mathcal P_g$, $g\in\mathcal B$, is
the set of primes in an interval contained in
$[\exp(L^c),\exp(L^d)]$.  In particular, intersecting a group with a
further interval again gives an interval of primes.  This additional
assumption allows us to use \Cref{thm:sw} with one label varying and all
other labels fixed.

\begin{theorem}[Construction of the residual ideal operator]
\label{thm:ideal}
Fix $E_{\rm id}>0$ and a fixed frequency exponent $C_4\geq0$.
There are fixed integers $m,J_0$ and a fixed constant $C_1$, depending
only on these parameters, $\ell$, and the prime-group bounds, with the
following property.  For every fixed $J\geq J_0$, there is a family
consisting of the raw pattern
\[
 C_g=\varnothing\quad\hbox{for every }g,\qquad K_0=1,
\]
and signed comparison patterns satisfying
\[
 C_g=\varnothing\ (g\in\mathcal S),\qquad
 |C_g|\leq m\ (g\in\mathcal B),\qquad
 \sum_\nu\sup|K_\nu|\leq L^{C_1},
\]
such that the symmetrized ideal operator of \Cref{sec:transference}
satisfies
\begin{equation}\label{eq:ideal-small-norm}
 \sup_{\Theta\in\R,\ |\zeta|\leq L^{C_4}}
 \norm{\mathcal T(\Theta)}_{2\to2}\leq L^{-E_{\rm id}}
\end{equation}
for sufficiently large $x$.  The coefficients and the family are
independent of $\Theta$ and of the individual value of $\zeta$.
Their defining parameters are independent of $J$; $J$ only supplies
the available shared slots.  The threshold for $x$ may depend on $J$.

More precisely, designate $m$ of the first $J$ slots in each big group
as probes, and split the big groups into two blocks. Each comparison
frees nonempty sets $A$ and $B$ of designated probes in the respective
blocks, and its coefficient has the form
\begin{equation}\label{eq:ideal-comparison-coefficient}
 -(-1)^{|A|+|B|}\mathcal K_A(D_A,X_A)
                         \mathcal K_B(D_B,X_B).
\end{equation}
Here $D_A,D_B$ are the products of the free labels used to form $D$,
$X_A$ is the product of the target labels in $A$, and $X_B$ the
product of the input labels in $B$.  These coefficients use no
shared labels and none of the last $\ell$ labels.  Each kernel is
a finite sum of products of matched log-cell indicators and pairs of
Dirichlet characters of conductor bounded by a fixed power of $L$.
\end{theorem}

The construction separates an approximation from a cancellation. We first
build kernels that replace selected shared products by independent products
in a bilinear pairing, for each prescribed fixed error exponent. The tests may depend on the entire ordered label tuples, as they
will when we condition on the labels outside the freed slots.

The cancellation then decomposes the input test in the first-block probes
and the target test in the second-block probes. Each decomposition is
orthogonal: its pieces are mean zero in a specified set of these probe
coordinates and independent of the others. Freeing $A$ averages the
source coordinates in $A$, and freeing $B$ averages the target coordinates
in $B$; the asymmetric coefficient above leaves these coordinates absent
from the multiplier. Thus only components independent of the freed
coordinates survive. Alternating subset signs cancel every pair for which the input component
is independent of at least one first-block probe and the target component
is independent of at least one second-block probe.
Every remaining pair contains a component that is mean zero in every
probe of one block. Such components have small norms by permutation
symmetry and the averaging of the last $\ell$ coordinates. We prove this cancellation after the
comparison estimate, so that its use on conditional tuple tests is explicit.

All label tuples in the next three subsections have their product
probability measure, with repetitions allowed.

\subsection{Products of labels and an elementary bilinear bound}

Let $\Lambda$ be a nonempty set of slots, using $k_g\leq m$ slots
from group $g$, and let $P_\Lambda$ denote the product of its labels.
Disjointness of the prime groups and unique factorization give, on
the support of this product,
\begin{equation}\label{eq:ideal-product-atom}
\begin{gathered}
 \rho_\Lambda(n):=\Prob(P_\Lambda=n)
 =\frac1n\prod_g
     \frac{k_g!}{V_g^{k_g}\prod_{p\in\mathcal P_g}a_p!}
 \leq\frac{L^\kappa}{n},\\
 \kappa=C_0\bigl(\log(m!)+m\log\max(1,v_-^{-1})\bigr),
\end{gathered}
\end{equation}
where $a_p$ is the multiplicity of $p$ in $n$.  The probability is
zero if the prescribed group multiplicities do not hold.  Indeed,
each factor apart from $1/n$ is at most
$m!\max(1,v_-^{-1})^m$, and there are at most $C_0T$ big groups.
In particular, this exponent does not depend on $J$.

This estimate applies to tests on the entire ordered label tuple.
If $f$ is any such $L^2$ test, supported where $c_1U\leq
P_\Lambda\leq c_2U$ for fixed $c_1,c_2>0$, put
$a_n=\E[f\ind_{P_\Lambda=n}]$.  Cauchy--Schwarz on each fiber gives
\begin{equation}\label{eq:ideal-collapse}
 \sum_n|a_n|^2
 \leq\sum_n\rho_\Lambda(n)
                 \E[|f|^2\ind_{P_\Lambda=n}]
 \leq\frac{L^\kappa}{c_1U}\norm f_2^2.
\end{equation}
No assumption that $f$ is a function of the product has been made.

\begin{lemma}[Integer bilinear estimate]\label{lem:bilinear}
Let $U,V\geq1$.  Suppose $a_n,b_t$ are supported in integer
intervals of lengths at most $C U,C V$, respectively, where $C$ is
fixed.  If
\[
 \alpha=u/r+\beta,\qquad (u,r)=1,\quad r\geq1,
 \qquad |\beta|\leq r^{-2},
\]
then
\begin{equation}\label{eq:integer-bilinear}
 \left|\sum_{n,t}a_nb_t\ee{\alpha nt}\right|
 \ll_C \norm a_{\ell^2}\norm b_{\ell^2}
 \left[U+(1+V/r)\{U+r\log(2r)\}\right]^{1/2}.
\end{equation}
The containing intervals need not begin at $1$, and the
coefficients may vanish on arbitrary subsets of them.
\end{lemma}

\begin{proof}
Extend the coefficients by zero to their containing intervals.
Cauchy--Schwarz in $n$, followed by expansion of the square and
the geometric-sum bound on the full interval of $n$, gives
\begin{align*}
 \left|\sum_{n,t}a_nb_t\ee{\alpha nt}\right|^2
 &\ll_C\norm a_{\ell^2}^2
 \left(U\sum_t|b_t|^2+
   \sum_{1\leq |h|\leq CV}
    \min\{U,\norm{\alpha h}_{\R/\Z}^{-1}\}
       \sum_t|b_t b_{t+h}|\right)\\
 &\ll_C\norm a_{\ell^2}^2\norm b_{\ell^2}^2
  \left(U+\sum_{1\leq h\leq CV}
         \min\{U,\norm{\alpha h}_{\R/\Z}^{-1}\}\right).
\end{align*}
At a zero denominator the minimum is interpreted as $U$.
For $r\geq2$, split the $h$-range into
$O_C(1+V/r)$ consecutive blocks, each of diameter at most $r/2$.
If $h\ne h'$ belong to a block, their difference is nonzero and
has magnitude less than $r$, so
\[
 \norm{\alpha(h-h')}_{\R/\Z}
 \geq\norm{u(h-h')/r}_{\R/\Z}
            -\frac{|h-h'|}{r^2}\geq\frac1{2r}.
\]
Thus the phases in a block are $1/(2r)$-separated on the circle.
Ordering them by distance from zero bounds their contribution by
\[
 O_C\left(U+\sum_{1\leq j\leq 2r}\frac rj\right)
 =O_C\bigl(U+r\log(2r)\bigr).
\]
For $r=1$, the trivial bound $O_C(UV)$ for the whole sum gives
the asserted estimate.  This proves the lemma.
\end{proof}

For two independent selected products of sizes $U,V$, combining
\Cref{eq:ideal-collapse,eq:integer-bilinear} shows that the
probability-space bilinear form with phase $\ee{\alpha P_A P_B}$
has norm at most
\begin{equation}\label{eq:ideal-probability-bilinear}
 C L^\kappa
 \left(\frac1V+\frac1r+\frac{\log(2r)}U+
                   \frac{r\log(2r)}{UV}\right)^{1/2}.
\end{equation}
The same assertion holds with the phase multiplied by
$(P_AP_B)^{i\zeta}$ for any real $\zeta$, since this factor
separates between the two tests and preserves their norms.

\subsection{Log cells and conditional character operators}

For a fixed nonempty slot set $\Lambda$, let $D,X$ be independent
products of label tuples of this type.  Given $Q\geq1$, $0<h\leq1$,
and $0<\xi\leq1$, put
\[
 I_j=[jh,(j+1)h),\qquad
 \nu_{\Lambda j}=\Prob(\log D\in I_j),\qquad
 \mathcal C_Q=\{\chi:\chi\text{ primitive},\ \operatorname{cond}\chi\leq Q\}.
\]
The conductor-$1$ character is included.  Define
\begin{equation}\label{eq:ideal-cell-kernel}
 \mathcal K_\Lambda(D,X)
 =\sum_{j:\nu_{\Lambda j}\geq\xi}
   \frac{\ind_{\log D\in I_j}\ind_{\log X\in I_j}}{\nu_{\Lambda j}}
        \sum_{\chi\in\mathcal C_Q}\overline{\chi(D)}\chi(X).
\end{equation}
We use $\mathcal K_A,\mathcal K_B$ for the corresponding slot sets.
All powers of $L$ chosen below are fixed before $J$.

There are at most $Q^2$ characters in $\mathcal C_Q$.  The number
$N_\Lambda$ of cells of positive mass satisfies
\begin{equation}\label{eq:ideal-cell-costs}
 N_\Lambda\ll_m1+T L^d/h,\qquad
 \sup|\mathcal K_\Lambda|\leq Q^2\xi^{-1},\qquad
 \sup_X\E_D|\mathcal K_\Lambda(D,X)|\leq Q^2.
\end{equation}
The last inequality follows because, for fixed retained $X$,
the mass $\nu_{\Lambda j}$ of its cell cancels the denominator.
The kernel is zero for discarded $X$ and has support only where
$|\log D-\log X|<h$.

We record explicitly the character estimate under a cell
restriction.  If $Q$ is a fixed power of $L$ and $\chi,\chi'$ are
distinct members of $\mathcal C_Q$, then, for every fixed $D_*>0$,
\begin{equation}\label{eq:ideal-character-cell}
 \left|\E\left[\ind_{\log D\in I_j}
          \overline{\chi(D)}\chi'(D)\right]\right|
 \ll_{D_*}L^{-D_*},
\end{equation}
uniformly in $j$ and in the nonempty slot set with at most $m$
slots per group.  To see this, hold all but one label fixed.
The cell and the selected prime group restrict the remaining
prime to an interval, possibly empty.  The character
$\overline\chi\chi'$ on the common modulus is nonprincipal:
otherwise the uniqueness of primitive induction would imply
$\chi=\chi'$.  Its modulus is at most $Q^2$.  For every fixed $K$,
 \Cref{thm:sw}, summed against the character over reduced
residue classes and applied with a larger accuracy exponent, gives
\[
 A_\rho(t):=\sum_{p\leq t}\rho(p)
       \ll_K t(\log t)^{-K}
\]
for these nonprincipal characters, uniformly for
$t\geq y:=\exp(L^c)$.  The allowed modulus is a fixed power of
$\log y$, as required there.  Partial summation on any subinterval
of $[y,\exp(L^d)]$ bounds its reciprocal-prime sum by
\[
 O_K\left((\log y)^{-K}+
            \int_y^\infty\frac{dt}{t(\log t)^K}\right)
 =O_K\bigl((\log y)^{1-K}\bigr).
\]
Divide by $V_g\geq v_-$ and choose $K$ sufficiently large in terms
of $D_*$.  Averaging the other labels proves
\Cref{eq:ideal-character-cell}; its precision is independent of
the width or location of the interval.

For a retained cell of mass $\nu$, the integral operator defined
by \Cref{eq:ideal-cell-kernel} is the character frame operator on
the conditional probability space:
\begin{equation}\label{eq:ideal-frame}
 f\longmapsto\sum_{\chi\in\mathcal C_Q}
       \chi\,\langle\chi,f\rangle_{L^2(\Prob(\cdot\mid I_j))}.
\end{equation}
Its nonzero eigenvalues are those of the character Gram matrix.
Every diagonal entry is $1$, since all selected primes exceed
$Q^2$ for large $x$.  By \Cref{eq:ideal-character-cell}, every
off-diagonal entry is $O_{D_*}(\xi^{-1}L^{-D_*})$.
The Gram matrix therefore has norm at most
\begin{equation}\label{eq:ideal-gram-bound}
 1+O_{D_*}(Q^2\xi^{-1}L^{-D_*})\leq2
\end{equation}
once the character precision has been chosen sufficiently large.
Passing from the conditional to the original measure multiplies
both squared norms by $\nu$, so the operator norm is unchanged.
Different cells are orthogonal blocks, and discarded cells have
zero operator.  Consequently kernel integration on the full
label space has norm at most $2$.  The adjoint and transpose have
the same bound, so the integration may be moved onto either test
in a bilinear form.  In particular neither a cell-count factor
nor a factor $\xi^{-1}$ is lost in this operator norm.

\subsection{Uniform comparison for arbitrary tuple tests}

\begin{lemma}[Comparison kernel]\label{lem:comparison-kernel}
Fix $m$, $A_0>0$, and $C_4\geq0$.  There are fixed powers
$Q,h^{-1},\xi^{-1}$ of $L$, independent of $J$, such that the
kernels in \Cref{eq:ideal-cell-kernel} have the following property.
Let $A,B$ be nonempty sets of slots, using at most $m$ slots per
big group.  Take independent label tuples for
$D_A,X_A,D_B,X_B$ with the prescribed slot laws.  If $f$ and $g$
are arbitrary $L^2$ tests on the tuples making up $X_B$ and
$X_A$, respectively, then for every $\alpha\in\R$ and
$|\zeta|\leq L^{C_4}$,
\begin{equation}\label{eq:ideal-comparison}
 \begin{split}
 \Big|\E\,\overline{f(X_B)}g(X_A)
  \big[&(X_AX_B)^{i\zeta}\ee{\alpha X_AX_B}\\
      &-\mathcal K_A(D_A,X_A)\mathcal K_B(D_B,X_B)
                (D_AD_B)^{i\zeta}\ee{\alpha D_AD_B}\big]\Big|
 \leq L^{-A_0}\norm f_2\norm g_2.
 \end{split}
\end{equation}
Here notation such as $f(X_B)$ denotes a tuple test, not an
assumption of dependence only on the product.  The kernels are
independent of $\alpha$ and $\zeta$.
\end{lemma}

\begin{proof}
Split the tests according to dyads $X_A\in[U,2U)$ and
$X_B\in[V,2V)$.  Each selected product lies between
$\exp(L^c)$ and $\exp(mC_0T L^d)$; thus for sufficiently large
$x$ there are at most $L^{d+1}$ dyads for each product.
We prove a bound $L^{-B_0}$ on each pair of dyads, with
$B_0>A_0+2(d+1)+10$.  Summing the bounds proves the lemma,
even using the crude bound of the original test norm for every
dyadic restriction.

Let $Z=UV$ and introduce a further fixed power $S$ of $L$.
Dirichlet approximation with maximum denominator
$\lfloor Z/S\rfloor$ gives a reduced fraction satisfying
\begin{equation}\label{eq:ideal-dirichlet}
 \alpha=u/r+\beta,\qquad
 1\leq r\leq Z/S,\qquad |\beta|\leq r^{-2},\qquad
 |\beta|Z\leq 2S/r.
\end{equation}
Here $Z/S\to\infty$ and the factor $2$ accounts for the integer
part.  This approximation is used only in the proof: the kernel
does not depend on the choice of the rational approximation.

\paragraph{Denominators larger than $Q$.}
Suppose $r>Q$.  Apply \Cref{eq:ideal-probability-bilinear} to
the raw term.  For the comparison term, move the two separate
kernel integrations onto the two tests.  Their norms increase
by at most $2$ each by \Cref{eq:ideal-gram-bound}.  The resulting
tests on $D_A,D_B$ are supported in $[U/2,4U]$ and $[V/2,4V]$
for large $x$, by the log localization.  The same bilinear bound
therefore applies with an absolute constant change.
In both uses we collapse arbitrary tuple tests by
\Cref{eq:ideal-collapse} before applying the integer lemma.

The two scales exceed every fixed power of $L$, and
$\log(2r)\leq L^{d+1}$ for large $x$.  Since $r>Q$ and
$r\leq UV/S$, the bound for either term, apart from an absolute
constant and the two test norms, is
\begin{equation}\label{eq:ideal-minor-error}
 L^\kappa
 \left(o(L^{-A})+Q^{-1}+L^{d+1}S^{-1}\right)^{1/2}
\end{equation}
for every fixed $A$.  Choosing $Q,S$ sufficiently large gives
the required dyadic precision.  This part is uniform in all
real $\zeta$, because the Mellin factor separates.

\paragraph{Denominators at most $Q$.}
Suppose $r\leq Q$.  All products under consideration are units
modulo $r$.  Fourier expansion on the finite group of units gives
\begin{equation}\label{eq:ideal-rational-expansion}
 \ee{uv/r}=\sum_{\psi\bmod r}c_\psi\psi(v),\qquad
 \sum_\psi|c_\psi|\leq\sqrt{\varphi(r)}\leq\sqrt Q.
\end{equation}
Indeed, in normalized counting measure the left side has
squared norm $1$, so Parseval gives $\sum|c_\psi|^2=1$ and
Cauchy--Schwarz proves the stated bound.  The unique primitive
character inducing $\psi$ has conductor at most $r$ and is
included in $\mathcal C_Q$.

First remove discarded cells.  Their total probability for
either product is at most
\begin{equation}\label{eq:ideal-discarded-mass}
 \Prob(\mathcal E)\leq(N_A+N_B)\xi,
 \qquad
 \mathcal E=\{X_A\text{ or }X_B\text{ lies in a discarded cell}\}.
\end{equation}
The comparison term is zero on this event.  The raw term there
has absolute value at most
\begin{align*}
 \E[\ind_{\mathcal E}|f(X_B)g(X_A)|]
 &\leq\Prob(\mathcal E)^{1/2}
       \bigl(\E|f(X_B)g(X_A)|^2\bigr)^{1/2}\\
 &=\Prob(\mathcal E)^{1/2}\norm f_2\norm g_2.
\end{align*}
The last equality uses independence of the two tuple spaces.
This argument also applies to tests concentrated in a discarded
cell; it does not assert that restriction to a small event is
small as an $L^2$ operator.

On retained cells, write the remaining scalar multiplier in the
log variable as
\[
 M(v)=\exp(i\zeta v)\ee{\beta e^v}.
\]
On the enlarged product range its derivative has magnitude
$O(L^{C_4}+S)$, by \Cref{eq:ideal-dirichlet}.  Since both
matched log products differ by less than $h$, we may replace
$M(\log(D_AD_B))$ by $M(\log(X_AX_B))$ while leaving the
rational phase unchanged.  For each fixed retained pair
$(X_A,X_B)$, \Cref{eq:ideal-cell-costs} bounds the error after
integration in $D_A,D_B$ by
\begin{equation}\label{eq:ideal-log-error}
 O\bigl(Q^4h(S+L^{C_4})\bigr).
\end{equation}
This is a pointwise error.  Pairing it with the tests costs at
most the same bound times $\norm f_2\norm g_2$, since their
spaces are independent.

For a character $\psi$ in \Cref{eq:ideal-rational-expansion},
the Gram calculation also gives the following pointwise
reproduction formula for retained $X$:
\begin{equation}\label{eq:ideal-character-reproduction}
 \E_D\mathcal K_\Lambda(D,X)\psi(D)
 =\psi(X)+O_{D_*}(Q^2\xi^{-1}L^{-D_*}).
\end{equation}
In fact, the summand corresponding to its inducing primitive
character is exactly $\psi(X)$: the diagonal conditional
expectation is $1$.  Each other summand is bounded by
\Cref{eq:ideal-character-cell} divided by the retained-cell
mass.  There are at most $Q^2$ such summands.  Choose $D_*$ so
that the displayed error is at most $1$.  Applying the formula
on both sides and then \Cref{eq:ideal-rational-expansion}
reproduces $\ee{uX_AX_B/r}$ with pointwise error
\begin{equation}\label{eq:ideal-major-character-error}
 O_{D_*}(Q^{5/2}\xi^{-1}L^{-D_*}).
\end{equation}
Multiplication by the fixed scalar $M(\log(X_AX_B))$ has
modulus one and does not change this error.  Together,
\Cref{eq:ideal-discarded-mass,eq:ideal-log-error,eq:ideal-major-character-error}
establish the desired comparison on a pair of dyads.

\paragraph{A noncircular choice of parameters.}
For completeness the following sufficient inequalities fix all
the precisions just used.  Write
\[
 Q=L^{q_*},\quad S=L^{s_*},\quad
 h=L^{-h_*},\quad \xi=L^{-x_*}.
\]
After $m,A_0,C_4$ have been fixed, choose successively
\begin{align}
 B_0&>A_0+2(d+1)+10,\notag\\
 q_*,s_*&>2(B_0+\kappa+10)+d+1,\notag\\
 h_*&>4q_*+\max(s_*,C_4)+B_0+10,\label{eq:ideal-parameter-order}\\
 x_*&>2B_0+d+1+h_*+10,\notag\\
 D_*&>B_0+x_*+\tfrac52q_*+10.\notag
\end{align}
Indeed $N_A,N_B\leq L^{d+1+h_*}$ for large $x$.
The successive lines control, respectively, dyadic summation,
\Cref{eq:ideal-minor-error}, \Cref{eq:ideal-log-error}, the
square root of \Cref{eq:ideal-discarded-mass}, and both
\Cref{eq:ideal-gram-bound,eq:ideal-major-character-error}.
Fixed implied constants are absorbed by the margins in these
inequalities.  Finally use \Cref{thm:sw} at the precision needed
for this $D_*$.  None of these choices involves $J$.
\end{proof}

\subsection{Symmetric probes and exact cancellation}

\begin{proof}[Proof of \Cref{thm:ideal}]
Split $\mathcal B$ into blocks $\mathcal B_1,\mathcal B_2$ of
sizes $s_1,s_2$ differing by at most one.  For sufficiently
large $x$ each size is at least $c_0T/3$.  In each group
designate the first $m$ of the first $J$ slots as probes; their
sets in the two blocks are $\mathcal I_1,\mathcal I_2$.
Let $P_i=|\mathcal I_i|=ms_i$ and $P=P_1+P_2$.
For every pair of nonempty subsets
$A\subseteq\mathcal I_1$, $B\subseteq\mathcal I_2$, free
exactly $A\cup B$ and use the coefficient in
\Cref{eq:ideal-comparison-coefficient}, with the kernels just
constructed.  The other comparison target slots, including
the last $\ell$, have exactly the independent laws in the
ideal operator's definition.

It suffices to bound pairings against symmetric vectors $F,H$,
since the operator has the orthogonal symmetrizing projection
on both sides.  In every pairing, average out the last $\ell$
input and target slots in each group.  No multiplier uses
these coordinates.  Denote the resulting functions on the
first $J$ coordinates by $F',H'$.  Both operations are
contractions.  The raw pairing is
\begin{equation}\label{eq:ideal-raw-pairing}
 R(F',H')=
 \E\left[\overline{F'}H'
       D_{\mathcal B}^{i\zeta}\ee{\Theta D_{\mathcal B}}\right],
\end{equation}
where input and target use the same first $J$ labels.  Its
bilinear norm is at most $1$.

We use the orthogonal product-space decomposition of
Efron and Stein~\cite[Section~2, Decomposition Lemma]{EfronStein1981},
and then count the components retained by the symmetry and probe
constraints. For a coordinate $i$, write $E_i$ for averaging its label.
Decompose $F'$ by the probes in the first block, and $H'$ by
those in the second:
\begin{align*}
 F'&=\sum_{S\subseteq\mathcal I_1}F_S,
 &F_S&=\prod_{i\in S}(1-E_i)
          \prod_{i\in\mathcal I_1\setminus S}E_iF',\\
 H'&=\sum_{T'\subseteq\mathcal I_2}H_{T'},
 &H_{T'}&=\prod_{i\in T'}(1-E_i)
          \prod_{i\in\mathcal I_2\setminus T'}E_iH'.
\end{align*}
These are orthogonal decompositions.  The indicated components
are separately mean zero in each indicated probe and
independent of the other probes of their block.  Put
$F_{\rm full}=F_{\mathcal I_1}$ and
$H_{\rm full}=H_{\mathcal I_2}$.

We need the precise effect of having first averaged out the
last $\ell$ slots.  In one group, decompose the original
symmetric vector on all $M=J+\ell$ coordinates by the same
coordinate averaging projections.  Among its components on
subsets of cardinality $j$, permutation symmetry makes their
squared norms equal.  Requiring all $m$ probes to occur in
the subset and all last $\ell$ slots to be absent retains
exactly the fraction
\begin{equation}\label{eq:ideal-hoeffding-fraction}
 \frac{\binom{M-m-\ell}{j-m}}{\binom Mj}
 =\frac{(j)_m(M-j)_\ell}{(M)_{m+\ell}}.
\end{equation}
The fraction is zero when a required cardinality is
impossible.  If $M\geq(m+\ell)(m+\ell-1)$, then
\[
 (M)_{m+\ell}
 =M^{m+\ell}\prod_{i=0}^{m+\ell-1}(1-i/M)
 \geq\tfrac12M^{m+\ell},
\]
using $\prod(1-a_i)\geq1-\sum a_i$ for $0\leq a_i\leq1$.
Consequently \Cref{eq:ideal-hoeffding-fraction} is at most
\begin{equation}\label{eq:ideal-hoeffding-bound}
 2(j/M)^m(1-j/M)^\ell
 \leq2\left(\frac m{m+\ell}\right)^m
          \left(\frac\ell{m+\ell}\right)^\ell
 \leq 2\ell^\ell m^{-\ell}=:\theta_m.
\end{equation}

This bound tensorizes even when the vectors are not products
across groups.  Indeed, first resolve the orthogonal
decomposition by the tuple of component cardinalities in
the groups of the block.  Independent within-group
permutations give equal squared norms for all tuples of
subsets with those cardinalities.  The fraction surviving
the stated inclusions and exclusions is the product of
\Cref{eq:ideal-hoeffding-fraction} over those groups.  Sum
the resulting inequality over the cardinality tuples.
Averaging the last $\ell$ coordinates in other groups is
a further contraction.  We obtain
\begin{equation}\label{eq:ideal-full-component}
 \norm{F_{\rm full}}_2\leq\theta_m^{s_1/2}\norm F_2,
 \qquad
 \norm{H_{\rm full}}_2\leq\theta_m^{s_2/2}\norm H_2.
\end{equation}
Choose $m$ so large that $\theta_m<1$ and
$-(c_0/6)\log\theta_m>E_{\rm id}+3$.
The right sides are then at most
$L^{-E_{\rm id}-3}\norm F_2$ and
$L^{-E_{\rm id}-3}\norm H_2$.  This is the only use of
large $m$, and it is possible because $\ell\geq1$.

Now fix a component pair $(F_S,H_{T'})$ and let
\[
 U=\mathcal I_1\setminus S,
 \qquad V=\mathcal I_2\setminus T'
\]
be its missing probes.  A comparison term with free sets
$A,B$ vanishes unless $A\subseteq U$ and $B\subseteq V$.
For if $i\in A\cap S$, the input label in coordinate $i$
is unshared, is absent from the multiplier, and is absent
from the target.  Integrating it annihilates the mean-zero
input component.  The same reasoning applies to a target
label in $B\cap T'$.  This uses the deliberate asymmetry in
\Cref{eq:ideal-comparison-coefficient}: its $A$ labels come
from the target, and its $B$ labels from the input.

Suppose the inclusions hold.  Condition on all shared
first-$J$ labels outside $A\cup B$, and write $C$ for their
product.  The input component is independent of the $A$
labels, and the target component is independent of the
$B$ labels.  Thus in the raw term the remaining tests are
an arbitrary tuple test on $X_B$ and one on $X_A$,
respectively, on independent product spaces.  In the
comparison term, the unused input $A$ labels and target
$B$ labels integrate out, leaving exactly the same two
tests and independent free products $D_A,D_B$.
The fixed shared product contributes $C^{i\zeta}$,
of modulus one, and changes the additive frequency to
$\alpha=\Theta C$.  \Cref{lem:comparison-kernel} therefore
matches the comparison term before its sign to
$R(F_S,H_{T'})$ with error at most
\begin{equation}\label{eq:ideal-component-error}
 L^{-A_0}\norm{F_S}_2\norm{H_{T'}}_2.
\end{equation}
To justify this bound after conditioning, apply the lemma
to the two conditional tests, average its error over the
shared labels, and use Cauchy--Schwarz.  The averages of
their squared conditional norms are precisely
$\norm{F_S}_2^2$ and $\norm{H_{T'}}_2^2$.
Uniformity in every real $\alpha$ is essential at this step.

When $U,V$ are nonempty, the signs give exact cancellation:
\begin{equation}\label{eq:ideal-inclusion-exclusion}
 \sum_{\substack{\varnothing\ne A\subseteq U\\
                  \varnothing\ne B\subseteq V}}
      (-1)^{|A|+|B|}
 =\left(\sum_{\varnothing\ne A\subseteq U}(-1)^{|A|}\right)
  \left(\sum_{\varnothing\ne B\subseteq V}(-1)^{|B|}\right)
 =(-1)(-1)=1.
\end{equation}
The initial minus sign in the comparison coefficient
therefore cancels the raw pairing for this component pair,
up to the errors just estimated.  If either missing set is
empty, no comparison survives.  All these exceptional raw
component pairs together are exactly
\[
 R(F_{\rm full},H')+
 R(F'-F_{\rm full},H_{\rm full}).
\]
We sum them before taking absolute values.  By
\Cref{eq:ideal-raw-pairing,eq:ideal-full-component}, their
total is at most
$2L^{-E_{\rm id}-3}\norm F_2\norm H_2$.
In particular there is no subset-count loss in this bound.

For the approximation errors, there are at most $3^P$
admissible triples consisting of a component pair and its
free subsets: each probe is either indicated, missing but
not freed, or freed.  Since $P\leq mC_0T$, choose
\[
 A_0>E_{\rm id}+mC_0\log3+10.
\]
Then the sum of \Cref{eq:ideal-component-error}, even
bounding each component norm by the original norm, is
at most $L^{-E_{\rm id}-10}\norm F_2\norm H_2$.
Together with the exceptional contribution this proves
\Cref{eq:ideal-small-norm} for sufficiently large $x$.

Finally, the number of comparisons is at most
$2^P\leq L^{mC_0\log2}$, and
\Cref{eq:ideal-cell-costs} bounds each coefficient by
$Q^4\xi^{-2}$.  For example, after making the choices in
\Cref{eq:ideal-parameter-order}, any fixed
\[
 C_1>mC_0\log2+4q_*+2x_*+1
\]
bounds the total coefficient cost, including the raw
pattern.  Expanding the two cell sums and two character
sums also has only a fixed power of $L$ terms.  Both free
products contain a prime, so $D_A,D_B\geq\exp(L^c)$.
The coefficient uses only the four products specified in
\Cref{eq:ideal-comparison-coefficient}; all assertions
about labels and subsequent separation follow directly
from \Cref{eq:ideal-cell-kernel}.

The order of choices is now explicit: choose $m$ from
\Cref{eq:ideal-full-component}, then $A_0$ to pay for
$3^P$, then the dyadic and kernel precisions in
\Cref{eq:ideal-parameter-order}, and hence $C_1$.
Only afterward require $J\geq m$ and
$J+\ell\geq(m+\ell)(m+\ell-1)$.  These requirements
define $J_0$.  A still larger fixed $J$ can therefore be
chosen to meet \Cref{thm:transference} without changing
any kernel parameter or the exponent $C_1$.
\end{proof}
